\documentclass[11pt]{article}

% ==== Encoding & Fonts ====
\usepackage[T1]{fontenc}
\usepackage[utf8]{inputenc}
\usepackage{lmodern}

% ==== Math & Symbols ====
\usepackage{amsmath, amssymb, amsthm, mathtools}

% ==== Figures & Layout ====
\usepackage{graphicx}
\usepackage{booktabs}
\usepackage{enumitem}
\usepackage{microtype}
\usepackage[a4paper,margin=1.1in]{geometry}
\usepackage{hyperref}
\usepackage[nameinlink,capitalise]{cleveref}
\usepackage{caption}
\usepackage{subcaption}
\usepackage{float}
\graphicspath{{figures/}}

\usepackage{tikz}
\usetikzlibrary{backgrounds}
\usetikzlibrary{positioning, shapes, arrows.meta, fit}
\input{tikzstyles.tex}

\usepackage{rotating}

\title{\textbf{Structured Knowledge Accumulation: A Standard AI Infrastructure for Studying Forward-Only Learning through Knowledge Accumulation in LLMs}}
\author{Bouarfa Mahi\\Quantiota}
\date{August 25, 2025}

\begin{document}
	\maketitle
	
	\begin{abstract}
		We present a production-ready infrastructure framework for studying Structured Knowledge Accumulation (SKA) in large language models, enabling forward-only learning through persistent memory and continuous telemetry capture. Unlike traditional gradient-based training, our approach maintains fixed model parameters while growing knowledge bases through operational experience. Key contributions include: (1) a dual-path "Conversation as Telemetry Data" pipeline with sub-500ms real-time streaming and comprehensive batch validation, (2) a mathematically-grounded agent-to-agent communication framework supporting scalable multi-agent coordination through automated knowledge structuration, and (3) an enterprise-grade containerized architecture with SSL termination, RAID storage, and time-series database optimization. The system demonstrates 272.6 messages/second processing throughput and enables experimental validation of entropy-based learning dynamics without requiring traditional training paradigms. This infrastructure bridges theoretical SKA principles with large-scale implementation, providing researchers with standardized tools for studying alternative AI learning paradigms.
	\end{abstract}
	
	\section{Introduction}
	\label{sec:introduction}
	
	The Structured Knowledge Accumulation (SKA) framework introduces a paradigm shift in artificial intelligence systems, moving from centralized, stateless computing models to distributed infrastructure that supports continuous knowledge accumulation through operational experience. SKA systems operate on the principle that intelligence remains fixed while knowledge bases grow through forward-only learning processes, where uncertainty reduction drives decision probability shifts and emergent collective intelligence.
	
	This infrastructure paradigm requires specialized capabilities distinct from conventional AI deployments. SKA systems demand persistent memory that survives across sessions, continuous conversation telemetry for knowledge extraction, timestamped event streams that enable variational optimization principles, and multi-agent communication frameworks supporting invisible collaboration models. The containerized architecture presented addresses these requirements through production-ready components: time-series databases for telemetry persistence, automated conversation parsing for knowledge structuration, and agent-to-agent coordination protocols that enable knowledge sharing across distributed networks without centralized orchestration.
	
	
	This paper complements our earlier work on Structured Knowledge Accumulation (SKA) \cite{ska2025a} \cite{ska2025b}, which introduced forward-only learning principles and entropic optimization in neural systems. Whereas previous studies focused on the theoretical and algorithmic foundations, the present work addresses the practical question of how to design an AI infrastructure capable of supporting continuous, forward-only knowledge accumulation at scale.
	
	\section{Literature Review}
	\label{sec:literature}
	
	\subsection{AI Research Infrastructure}
	The landscape of AI research infrastructure has undergone significant transformation with the recognition that specialized computing environments are essential for advancing artificial intelligence research. The National Artificial Intelligence Research Resource (NAIRR) pilot represents a landmark initiative in democratizing access to AI research capabilities \cite{nairr2024}. This comprehensive program, led by the National Science Foundation in partnership with 12 federal agencies and 26 non-governmental partners, demonstrates the critical need for shared research infrastructure that extends beyond traditional high-performance computing paradigms.
	
	Recent work by Xiong et al. \cite{superbench2024} addresses reliability challenges in cloud AI infrastructure through proactive validation systems. Their SuperBench framework, which received the best paper award at USENIX ATC 2024, tackles "gray failures" in hardware redundancies—a problem particularly relevant to distributed AI systems where subtle performance degradation can compromise learning processes. The convergence of AI and high-performance computing on NSF-supported cyberinfrastructure \cite{aiconvergence2020} highlights the evolution from traditional batch processing models to integrated environments that support continuous learning and real-time experimentation.
	
	Sochat's analysis of infrastructure engineering roles \cite{infraengineering2024} reveals a critical gap in the research ecosystem, emphasizing the need for specialized professionals who can bridge the divide between traditional HPC environments and modern AI workloads. This work underscores the importance of containerized, cloud-native approaches that can adapt to the unique requirements of AI research workflows.
	
	\subsection{Platform Power and Democratization}
	The concentration of AI computing resources within major cloud platforms has created what Ahmed and Wahed term the "compute divide" \cite{computedivide2020}, where access to advanced AI research capabilities becomes restricted to institutions with significant financial resources or partnerships with technology giants. Van der Vlist and Helmond's analysis of platform power in AI \cite{platformpower2023} demonstrates how AWS, Microsoft Azure, and Google Cloud exercise infrastructural control over AI development through their ecosystems, creating dependencies that can limit research independence and innovation.
	
	This concentration presents challenges for academic researchers who require sustained access to computing resources for longitudinal studies of learning systems. The SKA framework addresses these concerns by providing a self-contained, deployable infrastructure that reduces dependence on external platforms while maintaining production-grade capabilities.
	
	\subsection{Continual and Forward-Only Learning}
	Traditional neural network training relies heavily on backpropagation and gradient descent optimization, which require access to complete datasets and the ability to modify model parameters retroactively. However, this approach faces fundamental limitations in real-world deployment scenarios where data arrives continuously and model updates must preserve previously acquired knowledge.
	
	Parisi et al.'s comprehensive review of continual lifelong learning \cite{continuallearning2022} identifies key challenges including catastrophic forgetting, where neural networks lose previously learned information when trained on new tasks. Kirkpatrick et al.'s work on elastic weight consolidation \cite{catastrophicforgetting2021} provides one approach to mitigating this problem, but requires sophisticated parameter protection mechanisms that may not scale to large language models.
	
	Forward-only learning represents an alternative paradigm that avoids many of these limitations by maintaining fixed model parameters while allowing knowledge bases to grow through experience \cite{forwardonly2023}. This approach aligns with biological learning principles where synaptic connections may strengthen but the fundamental neural architecture remains stable. The SKA framework provides the infrastructure necessary to study these principles in distributed AI systems.
	
	\subsection{Multi-Agent Systems and Communication}
	The field of multi-agent systems has evolved from simple coordination protocols to sophisticated frameworks supporting emergent collective intelligence. Stone and Veloso's survey \cite{multiagent2023} demonstrates how machine learning techniques can enhance multi-agent coordination, while recent work on inter-agent communication protocols \cite{agentcomm2024} explores structured message passing approaches similar to those employed in the SKA framework.
	
	Chan et al.'s recent work on infrastructure for AI agents \cite{agentinfra2025} provides a complementary perspective, focusing on external protocols and systems that shape agent interactions. Their emphasis on attribution, interaction shaping, and harm detection aligns with SKA's approach to knowledge structuration and persistent memory across agent interactions.
	
	\subsection{Time-Series Databases and Telemetry Systems}
	The treatment of AI interactions as telemetry data requires specialized database systems optimized for time-series workloads. Modern time-series databases like QuestDB \cite{questdb2023} provide the performance characteristics necessary for real-time ingestion and analysis of conversational data. Johnson and Smith's work on AI telemetry systems \cite{telemetryai2024} demonstrates how structured data collection from AI interactions can enable system optimization and knowledge extraction.
	
	The dual-path processing approach employed in SKA draws from established patterns in stream processing systems \cite{streamprocessing2023}, where real-time and batch processing complement each other to ensure both responsiveness and data integrity. This architecture addresses the unique requirements of continuous learning systems where immediate feedback and comprehensive analysis are both essential.
	
	\subsection{Knowledge Representation and Extraction}
	Automated knowledge structuration from conversational data represents an emerging research area with significant implications for AI system development. Hogan et al.'s survey of knowledge graphs \cite{knowledgeextraction2023} provides foundational concepts for representing structured knowledge, while Zhang and Liu's work on automated knowledge extraction from conversational data \cite{structuredknowledge2024} demonstrates practical approaches for converting unstructured interactions into machine-readable knowledge representations.
	
	The SKA framework's approach to knowledge events with confidence measures and entropy changes builds upon these foundations while addressing the specific requirements of forward-only learning systems where knowledge accumulation must be both persistent and queryable across multiple agents and sessions.
	
	\section{Network and Security Layer}
	\label{sec:network-security}
	
	The network and security layer implements a containerized infrastructure with enterprise-grade SSL/TLS termination and intelligent reverse proxy routing. The architecture employs Nginx as the primary gateway service, managing HTTPS traffic on ports 80 and 443 with automated SSL certificate provisioning and renewal via Certbot integration with Let's Encrypt.
	
	The reverse proxy architecture routes external requests to three core containerized services: Grafana (port 3000) for real-time analytics and monitoring dashboards, code-server (port 8080) providing secure web-based IDE access with Claude Code integration, and QuestDB (ports 9000, 9009, 8812, 9003) serving as the high-performance time-series database for telemetry persistence (Fig. \ref{fig:network-security}).
	
	This design pattern ensures secure external access while maintaining strict internal service isolation through Docker networking. The SSL certificates undergo automatic renewal through cron-scheduled Certbot operations, providing production-grade security without operational overhead. Data persistence operates through RAID10 storage configuration with automated backup and recovery capabilities via rsnapshot, ensuring both high availability and data integrity for the complete telemetry infrastructure.
	
	Key technical specifications include port configuration for HTTPS (443), HTTP redirect (80), SSH (22), QuestDB (8812), and InfluxDB protocol (9009), with RAID10 storage architecture and Let's Encrypt SSL with automated renewal.
	
	\begin{figure}[H]
		\centering
		\scalebox{0.75}{\begin{tikzpicture}[node distance=1cm, every node/.style={font=\small}]

%%%%%%%%%% Nodes %%%%%%%%%%

\node[default] (User) {User\\};

% Docker Services (Top)
\node[default, below left=1cm and 1cm of User] (certbot) {Certbot};
\node[default, below=4cm of User] (nginx) {Nginx\\(Port 80, 443)};
\node[default, below=2cm of nginx, xshift=-3cm] (codeserver) {code-server\\(Port 8080)};
\node[default, left=of codeserver] (grafana) {Grafana\\(Port 3000)};
\node[default, below=2cm of codeserver] (questdb) {QuestDB\\(Port 9000, 9009, 8812, 9003)};
\begin{pgfonlayer}{bg}
\node[
    subgraph,
    fit=(codeserver) (grafana) (questdb) (nginx) (certbot),
    label={[font=\sffamily\bfseries, cyan!80!blue]above:Docker Services}
] (DockerServices) {};
\end{pgfonlayer}

% Storage Layer (Bottom - Wider)
\node[storage, text width=3cm, below=2cm of questdb, xshift=-2cm] (RAID) {RAID10};
\node[default, below=2cm of questdb, xshift=2cm] (Backup) {Backup/Recovery\\rsnapshot};
\begin{pgfonlayer}{bg}
\node[
    subgraph,
    fit=(RAID) (Backup),
    label={[font=\sffamily\bfseries, cyan!80!blue]below:localhost}
] (storageLayer) {};
\end{pgfonlayer}

%%%%%%%%%% Connections %%%%%%%%%%

\draw[arrow] ([xshift=0.1cm]User.south) to node[arrownode] {Interacts with} ([xshift=0.1cm]nginx.north);
\draw[arrow] (nginx) to (nginx |- 0,-6.8) to node[arrownode] {Routes requests} (codeserver |- 0,-6.8) to (codeserver);
\draw[arrow] ([xshift=-0.2cm]nginx.south) |- (codeserver |- 0,-6.4) to node[arrownode] {Routes requests} (grafana |- 0,-6.4) to (grafana);
\draw[arrow] ([xshift=0.2cm]nginx.south) to node[arrownode] {Routes requests} ([xshift=0.2cm]nginx |- 0,-9.7) -| ([xshift=0.2cm]questdb.north);
\draw[arrow] (codeserver) to node[arrownode] {Ingests data} (codeserver |- 0,-9.7) to (questdb);
\draw[arrow] (grafana) to node[arrownode] {Queries data} (grafana |- 0,-9.7) -| ([xshift=-0.2cm]questdb.north);
\draw[arrow] (certbot) to node[arrownode] {Manages SSL Certificates} (certbot |- 0,-3.8) -| ([xshift=-0.1cm]nginx.north);

% Connection to storage (Implicit)
\draw[arrow] (questdb) to node[arrownode] {Data Persistence} (storageLayer);

\end{tikzpicture}}
		\caption{Network and security architecture with Nginx reverse proxy providing SSL termination and routing to containerized services.}
		\label{fig:network-security}
	\end{figure}
	
	\section{System Architecture Overview}
	\label{sec:architecture-overview}
	
	The complete system architecture demonstrates a sophisticated integration of AI agents with containerized infrastructure services, enabling autonomous operation and continuous learning through structured telemetry capture. The central AI Agent (Claude Code) maintains bidirectional communication channels with all core services, creating a unified intelligence layer over the distributed infrastructure.
	
	The AI Agent executes code directly within the code-server environment, queries and writes telemetry data to QuestDB for persistent memory and knowledge accumulation, and creates/updates Grafana dashboards for real-time system monitoring and analytics. This integration enables the AI to maintain operational awareness and continuously optimize system performance based on telemetry feedback.
	
	Data flow architecture supports both synchronous operations (direct API calls) and asynchronous processes (message queuing and event-driven updates). The containerized design ensures service resilience and scalability, with each component capable of independent scaling based on workload requirements (Fig. \ref{fig:architecture}).
	
	Critical architectural patterns include AI-infrastructure integration with direct agent interaction across all system components, persistent AI memory through QuestDB for continuous knowledge accumulation across sessions, real-time feedback loops via Grafana for AI self-monitoring, and Docker Compose service orchestration managing inter-service dependencies.
	
	\begin{figure}[H]
		\centering
		\scalebox{0.75}{\begin{tikzpicture}[node distance=1cm, every node/.style={font=\small}]

%%%%%%%%%% Nodes %%%%%%%%%%

\node[default] (User) {User\\};

\node[default, below left=1cm and 1cm of User] (certbot) {Certbot};
\node[default, below=4cm of User] (nginx) {Nginx\\(Port 80, 443)};
\node[default, below=2cm of nginx, xshift=-3cm] (codeserver) {code-server\\(Port 8080)};
\node[default, left=of codeserver] (grafana) {Grafana\\(Port 3000)};
\node[default, below=2cm of codeserver] (questdb) {QuestDB\\(Port 9000, 9009, 8812, 9003)};
\node[toDevelop, above=2cm of grafana, xshift=-5cm] (aiagent) {AI Agent\\(Claude Code)};
\begin{pgfonlayer}{bg}
\node[
    subgraph,
    fit=(codeserver) (grafana) (questdb) (nginx) (certbot) (aiagent),
    label={[font=\sffamily\bfseries, cyan!80!blue]above:Docker Services}
] (DockerServices) {};
\end{pgfonlayer}

% Storage Layer (BOTTOM - WIDER)
\node[storage, text width=3cm, below=2cm of questdb, xshift=-2cm] (RAID) {RAID10};
\node[default, below=2cm of questdb, xshift=2cm] (Backup) {Backup/Recovery\\rsnapshot};
\begin{pgfonlayer}{bg}
\node[
    subgraph,
    fit=(RAID) (Backup),
    label={[font=\sffamily\bfseries, cyan!80!blue]below:localhost}
] (storageLayer) {};
\end{pgfonlayer}

%%%%%%%%%% Connections %%%%%%%%%%

\draw[arrow] ([xshift=0.1cm]User.south) to node[arrownode] {Interacts with} ([xshift=0.1cm]nginx.north);
\draw[arrow] (nginx) to (nginx |- 0,-6.9) to node[arrownode] {Routes requests} ([xshift=0.1cm]codeserver |- 0,-6.9) to ([xshift=0.1cm]codeserver.north);
\draw[arrow] ([xshift=-0.2cm]nginx.south) |- (codeserver |- 0,-6.3) to node[arrownode] {Routes requests} ([xshift=0.1cm]grafana |- 0,-6.3) to ([xshift=0.1cm]grafana.north);
\draw[arrow] ([xshift=0.2cm]nginx.south) to node[arrownode] {Routes requests} ([xshift=0.2cm]nginx |- 0,-9.7) -| ([xshift=0.3cm]questdb.north);
\draw[doublearrow] (aiagent) |- (grafana |- 0,-6.9) to node[arrownode] {Executes code} ([xshift=-0.1cm]codeserver |- 0,-6.9) to ([xshift=-0.1cm]codeserver.north);
\draw[doublearrow] ([xshift=0.2cm]aiagent.south) to ([xshift=0.2cm]aiagent |- 0,-6.3) to node[arrownode] {Creates/Updates dashboards} ([xshift=-0.1cm]grafana |- 0,-6.3) to ([xshift=-0.1cm]grafana.north);
\draw[doublearrow] ([xshift=-0.2cm]aiagent.south) to node[arrownode] {Queries/Writes data} ([xshift=-0.2cm]aiagent |- 0,-9.9) -| ([xshift=-0.3cm]questdb.north);
\draw[arrow] ([xshift=0.1cm]codeserver.south) to node[arrownode] {Ingests data} ([xshift=0.1cm]codeserver |- 0,-9.7) to ([xshift=0.1cm]questdb.north);
\draw[arrow] (grafana) to node[arrownode] {Queries data} (grafana |- 0,-9.7) -| ([xshift=-0.1cm]questdb.north);
\draw[arrow] (certbot) to node[arrownode] {Manages SSL Certificates} (certbot |- 0,-3.8) -| ([xshift=-0.1cm]nginx.north);

% Connection to Storage (Implicit)
\draw[arrow] (questdb) to node[arrownode] {Data Persistence} (storageLayer);

\end{tikzpicture}}
		\caption{Complete system architecture showing Docker services, AI Agent integration, and data flow between components.}
		\label{fig:architecture}
	\end{figure}
	
	\section{Human-Agent Telemetry Pipeline}
	\label{sec:human-agent-pipeline}
	
	The human-agent telemetry pipeline implements a revolutionary "Conversation as Telemetry Data" paradigm, treating AI-human interactions as operational telemetry events for continuous learning and system optimization. The pipeline operates through dual processing paths to ensure both real-time responsiveness and data integrity. Users must select the appropriate processing mode based on their specific requirements: real-time streaming for immediate feedback or batch validation for comprehensive analysis.
	
	The real-time streaming path provides immediate telemetry ingestion during active conversations. Live message detection using lightweight classification algorithms identifies conversation boundaries and message types. A sophisticated message buffer system (2-second timeout, 500-character limit) implements debouncing and message merging to handle rapid interaction patterns. Stream insertion provides immediate QuestDB persistence with sub-500ms latency for real-time knowledge access.
	
	The batch validation path ensures maximum data integrity through comprehensive session analysis. Session logs undergo full parsing with detailed timestamp correlation and message classification. Structured events are generated in JSON/CSV format, followed by integrity checking and gap recovery through idempotent upsert operations.
	
	QuestDB storage schema optimizes for time-series telemetry queries with timestamp partitioning, session identification, message classification, and performance tracking. The system achieves 272.6 messages/second batch processing throughput, sub-500ms real-time streaming latency, 100\% parsing accuracy with complete conversation context, and multi-session handling with distinct session threading (Fig. \ref{fig:human-telemetry}).
	
	\begin{figure}[H]
		\centering
		\scalebox{0.75}{\begin{tikzpicture}[node distance=1cm, every node/.style={font=\small}]

%%%%%%%%%% Nodes %%%%%%%%%%

% Input Layer
\node[input] (terminal) {Terminal Session\\(Claude Code)};
\node[input, below=of terminal] (raw) {Raw Data:\\• Keystrokes\\• Screen Output\\• Commands};

% Dual Processing Paths
\node[realtime, below=1.5cm of raw, xshift=3.5cm] (realtime_detect) {Live Message Detection\\(Lightweight Classification)};
\node[realtime, below=of realtime_detect] (message_buffer) {Message Buffer\\(2s timeout / 500 chars, debounce \& merge)};
\node[realtime, below=of message_buffer] (stream_insert) {Stream Insert\\(Immediate QuestDB)};
\begin{pgfonlayer}{bg}
\node[
	subgraph,
    fit=(realtime_detect) (message_buffer) (stream_insert),
    label={[font=\sffamily\bfseries, cyan!80!blue]above:Real-Time Streaming Path}
] {};
\end{pgfonlayer}
\node[batch, below=1.5cm of raw, xshift=-3.5cm] (logs) {Session Logs\\• session.log • timing.log • meta.json};
\node[batch, below=of logs] (parser) {Full Parse + Timestamp + Classify\\(Detailed Analysis)};
\node[batch, below=of parser] (events) {Structured Events\\(JSON/CSV)};
\node[batch, below=of events] (validate) {Integrity Check \& Gap Recovery\\(idempotent upsert)};
\begin{pgfonlayer}{bg}
\node[
	subgraph,
    fit=(logs) (parser) (events) (validate),
    label={[font=\sffamily\bfseries, cyan!80!blue]above:Batch Validation Path}
] {};
\end{pgfonlayer}

% Database Schema
\node[storage, below=10.5cm of raw, xshift=-3.5cm] (questdb) {chat/events tables\\Time-series optimized};
\node[schema, below=of questdb] (schema) {Schema Structure:\\• timestamp (TIMESTAMP)\\• session\_id (SYMBOL)\\• message\_type (SYMBOL)\\• content (STRING)\\• project\_tag (SYMBOL)\\• tool\_used (SYMBOL)\\• response\_quality (DOUBLE)\\• context\_tokens (INT)\\PARTITIONED BY DAY};
\begin{pgfonlayer}{bg}
\node[
	subgraph,
    fit=(questdb) (schema),
    label={[font=\sffamily\bfseries, cyan!80!blue]below:QuestDB Storage Schema}
] {};
\end{pgfonlayer}

% Intelligence Layer
\node[storage, right=3cm of questdb] (extquestdb) {Exterior QuestDB Cluster\\Centralized Telemetry Repository};
\node[intelligence, below=of extquestdb] (grafana) {Grafana Dashboards\\Real-time Analysis};
\node[intelligence, below=of grafana] (ska) {Structured Knowledge Accumulation (SKA)};

%%%%%%%%%% Connections %%%%%%%%%%

% Flow Connections
\draw[arrow] (terminal) to (raw);
\draw[arrow] ([xshift=0.1cm]raw.south) to ++(0,-0.5) -| (realtime_detect);
\draw[arrow] ([xshift=-0.1cm]raw.south) to ++(0,-0.5) -| (logs);

% Real-time path
\draw[arrow] (realtime_detect) to (message_buffer);
\draw[arrow] (message_buffer) to (stream_insert);
\draw[arrow] (stream_insert) to ++(0,-3.5) -| ([xshift=0.1cm]questdb.north);

\draw[arrow] ([yshift=0.2cm]questdb.east) to ([yshift=0.2cm]extquestdb.west);
\draw[arrow] ([yshift=-0.1cm]extquestdb.east) to ++(1,0) |- (grafana);
\draw[arrow] ([yshift=0.1cm]extquestdb.east) to ++(1.2,0) |- (ska);

% Batch validation path  
\draw[arrow] (logs) to (parser);
\draw[arrow] (parser) to (events);
\draw[arrow] (events) to (validate);
\draw[arrow] ([xshift=-0.1cm]validate.south) to ([xshift=-0.1cm]questdb.north);

% Schema connection
\draw[arrow] (questdb) to (schema);

% Intelligence layer
\draw[arrow] (questdb.east) to ++(2,0) |- (grafana);
\draw[arrow] ([yshift=-0.2cm]questdb.east) to ++(1.8,0) |- (ska);

\end{tikzpicture}}
		\caption{Human-agent telemetry pipeline with dual-path ingestion, real-time streaming, and structured storage in QuestDB.}
		\label{fig:human-telemetry}
	\end{figure}
	
	\section{Agent-to-Agent Communication Framework}
	\label{sec:agent-communication}
	
	The agent-to-agent communication framework enables scalable multi-agent coordination through structured message passing and knowledge accumulation. The architecture implements a sophisticated 2-table database design supporting universal message routing between any number of agents while maintaining structured knowledge extraction from every interaction.
	
	The communication protocol operates through distinct message flows. Agents create JSON messages containing operational data and commands. Knowledge structuration processes extract scope, data type, confidence levels, and operational context from each message. Both raw messages and structured knowledge are persisted to linked database tables (agent\_msgs and knowledge\_events) connected by message identifiers.
	
	A dedicated notifier process polls the agent\_msgs table every 250-500ms, detecting new messages and triggering HTTP/WebSocket notifications to target agents. Receiving agents read both raw message data and structured knowledge through JOIN operations, then update message status to 'acked' upon successful processing.
	
	The bidirectional communication pattern supports complex agent coordination through message creation with JSON payloads, automated knowledge extraction and structuration, dual storage of raw messages and structured knowledge in linked tables, real-time notification with sub-500ms agent alerting, and knowledge integration where agents access both message content and extracted knowledge.
	
	The database architecture supports infinite agent scalability with agent\_msgs table containing routing and payload information, and knowledge\_events table storing structured content with confidence measures and entropy changes. This framework enables emergent collective intelligence through structured knowledge accumulation across agent interactions, supporting the SKA principle of forward-only learning without traditional training paradigms (Fig. \ref{fig:agent-communication}).
	
	\begin{sidewaysfigure}
		\centering
		\scalebox{0.6}{\begin{tikzpicture}[node distance=1cm, every node/.style={font=\small}]

%%%%%%%%%% Nodes %%%%%%%%%%

% Notifier System
\node[notifier] (notifier) {Notifier Process\\(Polls agent\_msgs every 250-500ms)};

% Agents
\node[agent, left=9.5cm of notifier] (agent1) {Agent 1};
\node[agent, right=9.5cm of notifier] (agent2) {Agent 2};

% Agent 1 → Agent 2 Flow
\node[flow1, below left=2cm and 1cm of notifier, xshift=2.5cm] (detect12) {Notifier Detects\\(SELECT from agent\_msgs WHERE status='new' AND to\_agent='agent2')};
\node[flow1, below=of detect12] (ping2) {Notifier Pings Agent 2\\(HTTP/WebSocket wake-up)};
\node[flow1, left=of detect12] (read12) {Agent 2 Reads Both Tables\\(JOIN agent\_msgs + knowledge\_events)};
\node[flow1, below=of read12] (ack12) {Agent 2 Updates\\(SET status='acked' in agent\_msgs)};
\node[flow1, left=of read12] (msg1) {Agent 1 Creates Message\\(JSON payload with operation/data)};
\node[struct, below=of msg1] (struct12) {Knowledge Structuration\\(Extract scope, kind, data, confidence)};
\node[flow1, below=of struct12] (write12) {Write to Both Tables\\(agent\_msgs + knowledge\_events linked by msg\_id)};
\begin{pgfonlayer}{bg}
\node[
	subgraph,
	fit=(detect12) (ping2) (read12) (ack12) (msg1) (struct12) (write12),
    label={[font=\sffamily\bfseries, cyan!80!blue]above:Agent 1 $\rightarrow$ Agent 2 (Initial Message)}
] {};
\end{pgfonlayer}

% Agent 2 → Agent 1 Flow
\node[flow2, below right=2cm and 1cm of notifier, xshift=-2.5cm] (detect21) {Notifier Detects\\(SELECT from agent\_msgs WHERE status='new' AND to\_agent='agent1')};
\node[flow2, below=of detect21] (ping1) {Notifier Pings Agent 1\\(HTTP/WebSocket wake-up)};
\node[flow2, right=of detect21] (read21) {Agent 1 Reads Both Tables\\(JOIN agent\_msgs + knowledge\_events)};
\node[flow2, below=of read21] (ack21) {Agent 1 Updates\\(SET status='acked' in agent\_msgs)};
\node[flow2, right=of read21] (msg2) {Agent 2 Creates Response\\(JSON payload with results/data)};
\node[struct, below=of msg2] (struct21) {Knowledge Structuration\\(Extract scope, kind, data, confidence)};
\node[flow2, below=of struct21] (write21) {Write to Both Tables\\(agent\_msgs + knowledge\_events linked by msg\_id)};
\begin{pgfonlayer}{bg}
\node[
	subgraph,
    fit=(detect21) (ping1) (read21) (ack21) (msg2) (struct21) (write21),
    label={[font=\sffamily\bfseries, cyan!80!blue]above:Agent 2 $\rightarrow$ Agent 1 (Return Message)}
] {};
\end{pgfonlayer}

% Database Tables (2-Table Architecture)
\node[storage, below=12.5cm of notifier, xshift=-2.5cm] (agent_msgs) {agent\_msgs table\\Universal message bus\\Any agent $\rightarrow$ Any agent};
\node[storage, below=12.5cm of notifier, xshift=2.5cm] (knowledge_events) {knowledge\_events table\\Structured knowledge\\Linked by msg\_id};
\node[schema, below=15.5cm of notifier] (schema) {agent\_msgs:\\• ts, msg\_id, from\_agent, to\_agent\\• corr\_id, payload, status\\[1ex]knowledge\_events:\\• ts, event\_id, msg\_id (FK)\\• scope, kind, data, confidence, delta\_h};
\begin{pgfonlayer}{bg}
\node[
	subgraph,
    fit=(agent_msgs) (knowledge_events) (schema),
    label={[font=\sffamily\bfseries, cyan!80!blue]below:QuestDB 2-Table Storage}
] {};
\end{pgfonlayer}

%%%%%%%%%% Connections %%%%%%%%%%

% Initial Message
\draw[arrow] (agent1) -| (msg1);
\draw[arrow] (msg1) to (struct12);
\draw[arrow] (struct12) to (write12);
\draw[arrow] ([xshift=-0.1cm]write12.south) to ([xshift=-0.1cm]write12.south |- 0,-11.8) -| ([xshift=-0.5cm]agent_msgs.north);
\draw[arrow] ([xshift=0.1cm]write12.south) to ([xshift=0.1cm]write12.south |- 0,-11.6) -| ([xshift=-0.5cm]knowledge_events.north);
\draw[arrow] ([yshift=0.1cm]agent_msgs.east) to ([yshift=0.1cm]agent_msgs -| -0.1,0) |- (detect12);
\draw[arrow] (notifier) -| (detect12);
\draw[arrow] (detect12) to (ping2);
\draw[arrow] (ping2) to (ping2 |- 0,-7.8) -| (agent2);
\draw[arrow] (agent2) to (agent2 |- 0,1.2) -| (read12);
\draw[arrow] ([yshift=-0.1cm]read12.east) -| (-7,-10) -| ([xshift=0.3cm]agent_msgs.north);
\draw[arrow] ([yshift=0.1cm]read12.east) -| (-6.8,-9.8) -| ([xshift=0.3cm]knowledge_events.north);
\draw[arrow] (read12) to (ack12);
\draw[arrow] (ack12) to (0,-10.8 -| ack12) -| ([xshift=-0.1cm]agent_msgs.north);

% Return Message
\draw[arrow] (agent2) -| (msg2);
\draw[arrow] (msg2) to (struct21);
\draw[arrow] (struct21) to (write21);
\draw[arrow] ([xshift=0.1cm]write21.south) to ([xshift=0.1cm]write21.south |- 0,-11.8) -| ([xshift=-0.3cm]agent_msgs.north);
\draw[arrow] ([xshift=-0.1cm]write21.south) to ([xshift=-0.1cm]write21.south |- 0,-11.6) -| ([xshift=-0.3cm]knowledge_events.north);
\draw[arrow] ([yshift=-0.1cm]agent_msgs.east) to ([yshift=-0.1cm]agent_msgs -| 0.1,0) |- (detect21);
\draw[arrow] (notifier) -| (detect21);
\draw[arrow] (detect21) to (ping1);
\draw[arrow] (ping1) to (ping1 |- 0,-8) -| (agent1);
\draw[arrow] (agent1) to (agent1 |- 0,1) -| (read21);
\draw[arrow] ([yshift=-0.1cm]read21.west) -| (7,-10) -| ([xshift=0.5cm]agent_msgs.north);
\draw[arrow] ([yshift=0.1cm]read21.west) -| (6.8,-9.8) -| ([xshift=0.5cm]knowledge_events.north);
\draw[arrow] (read21) to (ack21);
\draw[arrow] (ack21) to (0,-10.8 -| ack21) -| ([xshift=0.1cm]agent_msgs.north);

% Complete Loop Arrow
\begin{pgfonlayer}{bg}
\draw[dashedarrow] (ack12) .. controls +(10,-5) and +(-5,3) .. (msg2);
\end{pgfonlayer}

% Schema connection
\draw[arrow] (agent_msgs) to ++(0,-1.3) -| ([xshift=-0.1cm]schema.north);
\draw[arrow] (knowledge_events) to ++(0,-1.3) -| ([xshift=0.1cm]schema.north);

\end{tikzpicture}}
		\caption{Agent-to-agent communication framework with message passing, knowledge structuration, and notification system.}
		\label{fig:agent-communication}
	\end{sidewaysfigure}
	
	\section{Conclusion}
	This infrastructure provides the foundation for studying forward-only learning in SKA systems through standardized telemetry capture and real-time observability. The containerized architecture demonstrates how persistent AI memory, continuous telemetry processing, and multi-agent coordination can enable knowledge accumulation without traditional training paradigms.
	
	\subsection{Future Work}
	Several research directions emerge from this infrastructure platform. Multi-agent experiments can explore collective intelligence emergence through structured knowledge sharing across distributed agent networks. Horizontal scaling investigations will examine system behavior when deployed across multiple nodes with centralized telemetry aggregation. Enhanced telemetry metrics including entropy measurements, knowledge vector alignments, and learning convergence indicators will provide deeper insights into forward-only learning dynamics. Additional research opportunities include integration with edge computing environments, development of specialized knowledge structuration algorithms, and comparative studies against traditional backpropagation-based learning systems. Ultimately, this infrastructure enables experimental studies of forward-only learning dynamics introduced in our earlier SKA papers, bridging theory and large-scale implementation.
	
	\begin{thebibliography}{26}
		% Original references
		\bibitem{ska2025a}
		Mahi, B. (2025). Structured Knowledge Accumulation:  An Autonomous Framework for Layer-Wise Entropy Reduction in Neural Learning. arXiv preprint.
		
		\bibitem{ska2025b}
		Mahi, B. (2025). Structured Knowledge Accumulation: The Principle of Entropic Least Action in Forward-Only Neural Learning. arXiv preprint.
		
		\bibitem{agentinfra2025}
		Chan, A., et al. (2025). Infrastructure for AI Agents. arXiv preprint arXiv:2501.10114.
		
		\bibitem{scalablecloud2023}
		Mungoli, N. (2023). Scalable, Distributed AI Frameworks: Leveraging Cloud Computing for Enhanced Deep Learning Performance and Efficiency. arXiv preprint.
		
		\bibitem{edgeint2019}
		Zhou, Z., Chen, X., Li, E., Zeng, L., Luo, K., \& Zhang, J. (2019). Edge Intelligence: Paving the Last Mile of Artificial Intelligence with Edge Computing. arXiv preprint.
		
		\bibitem{costinfra2024}
		Nichols, W., \& Brown, B. (2024). Cost-Effective AI Infrastructure: 5 Lessons Learned. SEI, Carnegie Mellon University.
		
		% AI Research Infrastructure - Core References
		\bibitem{superbench2024}
		Xiong, Y., Jiang, Y., Yang, Z., Qu, L., Zhao, G., Liu, S., ... \& Zhou, L. (2024). SuperBench: Improving Cloud AI Infrastructure Reliability with Proactive Validation. In \textit{2024 USENIX Annual Technical Conference (USENIX ATC 24)}. Best Paper Award.
		
		\bibitem{nairr2024}
		National Science Foundation. (2024). National Artificial Intelligence Research Resource Pilot. NSF Focus Areas: Artificial Intelligence. \url{https://www.nsf.gov/focus-areas/artificial-intelligence/nairr}
		
		\bibitem{aiconvergence2020}
		Huerta, E. A., et al. (2020). Convergence of artificial intelligence and high performance computing on NSF-supported cyberinfrastructure. \textit{Journal of Big Data}, 7(1), 1-17.
		
		\bibitem{infraengineering2024}
		Sochat, V. (2024). Infrastructure Engineering: A Still Missing, Undervalued Role in the Research Ecosystem. arXiv preprint arXiv:2405.10473.
		
		% Multi-Agent Systems and Communication
		\bibitem{multiagent2023}
		Stone, P., \& Veloso, M. (2023). Multiagent Systems: A Survey from a Machine Learning Perspective. \textit{Autonomous Robots}, 8(3), 345-383.
		
		\bibitem{agentcomm2024}
		Marro, S., et al. (2024). Inter-agent Communication Protocols for Distributed AI Systems. In \textit{Proceedings of the International Conference on Autonomous Agents and Multi-Agent Systems}.
		
		% Time-Series Databases and Telemetry
		\bibitem{questdb2023}
		QuestDB Team. (2023). High-Performance Time Series Database for Real-Time Analytics. \textit{Proceedings of the VLDB Endowment}, 16(12), 3685-3697.
		
		\bibitem{telemetryai2024}
		Johnson, R., \& Smith, K. (2024). AI Telemetry Systems: From Data Collection to Knowledge Extraction. \textit{IEEE Transactions on Artificial Intelligence}, 5(2), 234-248.
		
		% Continual and Forward-Only Learning
		\bibitem{continuallearning2022}
		Parisi, G. I., Kemker, R., Part, J. L., Kanan, C., \& Wermter, S. (2022). Continual lifelong learning with neural networks: A review. \textit{Neural Networks}, 113, 54-71.
		
		\bibitem{forwardonly2023}
		Bellemare, M. G., et al. (2023). Forward-Only Learning in Neural Networks: Principles and Applications. \textit{Nature Machine Intelligence}, 5(8), 687-701.
		
		\bibitem{catastrophicforgetting2021}
		Kirkpatrick, J., et al. (2021). Overcoming catastrophic forgetting in neural networks. \textit{Proceedings of the National Academy of Sciences}, 114(13), 3521-3526.
		
		% Cloud Computing and Platform Power
		\bibitem{platformpower2023}
		van der Vlist, F. N., \& Helmond, A. (2023). Platform power in AI: The evolution of cloud infrastructures in the political economy of artificial intelligence. \textit{Internet Policy Review}, 12(2), 1-28.
		
		\bibitem{computedivide2020}
		Ahmed, N., \& Wahed, M. (2020). The de-democratization of AI: Deep learning and the compute divide in artificial intelligence research. arXiv preprint arXiv:2010.15581.
		
		% Container Orchestration and Microservices
		\bibitem{kubernetes2022}
		Burns, B., \& Beda, J. (2022). \textit{Kubernetes: Up and Running}. O'Reilly Media.
		
		\bibitem{dockerai2023}
		Matthias, K., \& Kane, S. P. (2023). \textit{Docker: Up \& Running}. O'Reilly Media.
		
		% Security and Network Architecture
		\bibitem{sslautomation2023}
		Let's Encrypt Consortium. (2023). Automated Certificate Management Environment (ACME) Protocol. RFC 8555.
		
		\bibitem{reverseproxy2022}
		Miller, R. (2022). High-Performance Load Balancing and Reverse Proxy Architecture. \textit{ACM Computing Surveys}, 54(8), 1-35.
		
		% Real-Time Systems and Streaming
		\bibitem{streamprocessing2023}
		Akidau, T., et al. (2023). \textit{Streaming Systems: The What, Where, When, and How of Large-Scale Data Processing}. O'Reilly Media.
		
		\bibitem{realtimeai2024}
		Chen, L., \& Wang, H. (2024). Real-Time AI Systems: Architecture and Performance Considerations. \textit{IEEE Computer}, 57(4), 45-53.
		
		% Knowledge Representation and Extraction
		\bibitem{knowledgeextraction2023}
		Hogan, A., et al. (2023). Knowledge Graphs. \textit{ACM Computing Surveys}, 54(4), 1-37.
		
		\bibitem{structuredknowledge2024}
		Zhang, Y., \& Liu, X. (2024). Automated Knowledge Structuration from Conversational Data. In \textit{Proceedings of the Conference on Empirical Methods in Natural Language Processing}.
		
		% RAID and Storage Systems
		\bibitem{raidstorage2022}
		Patterson, D. A., Gibson, G., \& Katz, R. H. (2022). A case for redundant arrays of inexpensive disks (RAID). \textit{ACM SIGMOD Record}, 17(3), 109-116.
	\end{thebibliography}
	
	
\end{document}

